\begin{pastexam}[label={px:gk3-115-ds-04}]{高考三級資訊處理 115 年 第四題}
（申論；25 分）

考慮以下兩個互相呼叫（mutually recursive）的 C 語言函式：
\begin{minted}[fontsize=\small]{c}
int foo(int n) {
    if (n <= 1) return 1;
    return foo(n - 1) + bar(n - 1) + 2;
}
int bar(int n) {
    if (n <= 1) return 1;
    return foo(n - 1) + bar(n - 1);
}
\end{minted}
請回答下列問題：（每一小題請寫出推導過程，無推導過程不予計分。）
\begin{itemize}[leftmargin=2.2em]
\item[㈠] 當執行 foo(10)時，一共會呼叫 foo()函式幾次（包含最外層 foo(10)的這一次呼叫）？（10 分）
\item[㈡] 執行 foo(10)的最終回傳結果為何？（10 分）
\item[㈢] 以 Big-O 表示 foo($n$)的時間複雜度（time complexity）。本題若同一個「函式與參數」組合被呼叫多次，每次
都重新計算，不會儲存先前的計算結果供之後使用。（5 分）
\end{itemize}
\end{pastexam}

\begin{pastsoln}
答案：㈠ $512$ 次。 ㈡ $1024$。 ㈢ $O(2^n)$。

\textbf{㈠} 令 $F(n)$ 是執行 foo($n$) 時 foo 被呼叫的總次數（含這一次），$B(n)$ 是執行 bar($n$) 時 foo 被呼叫的
總次數。 $n \le 1$ 時兩個函式都直接回傳，不再呼叫別人：$F(1) = 1$、$B(1) = 0$。 $n \ge 2$ 時，foo($n$) 自己算
$1$ 次，再呼叫 foo($n-1$) 和 bar($n-1$)；bar($n$) 自己不算 foo，再呼叫 foo($n-1$) 和 bar($n-1$)：
\[
F(n) = 1 + F(n-1) + B(n-1), \qquad B(n) = F(n-1) + B(n-1) \qquad (n \ge 2).
\]
兩式相減得 $F(n) = B(n) + 1$（$n = 1$ 時也成立）。 代回第二式：
\[
B(n) = \bigl(B(n-1) + 1\bigr) + B(n-1) = 2B(n-1) + 1,
\]
所以 $B(n) + 1 = 2\bigl(B(n-1) + 1\bigr)$，$B(n) + 1 = 2^{n-1}\bigl(B(1) + 1\bigr) = 2^{n-1}$，
$F(n) = B(n) + 1 = 2^{n-1}$。 逐項列出：
\begin{center}
\small
\begin{tabular}{l*{10}{c}}
\toprule
$n$ & 1 & 2 & 3 & 4 & 5 & 6 & 7 & 8 & 9 & 10 \\
\midrule
$F(n)$ & 1 & 2 & 4 & 8 & 16 & 32 & 64 & 128 & 256 & 512 \\
$B(n)$ & 0 & 1 & 3 & 7 & 15 & 31 & 63 & 127 & 255 & 511 \\
\bottomrule
\end{tabular}
\end{center}
執行 foo(10) 共呼叫 foo \textbf{$512$ 次}（bar 被呼叫 $511$ 次）。

\textbf{㈡} 令 $f(n)$、$b(n)$ 是 foo($n$)、bar($n$) 的回傳值：$f(1) = b(1) = 1$，$n \ge 2$ 時
\[
f(n) = f(n-1) + b(n-1) + 2, \qquad b(n) = f(n-1) + b(n-1).
\]
令 $s(n) = f(n) + b(n)$，則 $n \ge 2$ 時 $f(n) = s(n-1) + 2$、$b(n) = s(n-1)$，相加得
\[
s(n) = 2s(n-1) + 2, \qquad s(1) = 2 .
\]
兩邊加 $2$：$s(n) + 2 = 2\bigl(s(n-1) + 2\bigr)$，所以 $s(n) + 2 = 2^{n-1}\bigl(s(1) + 2\bigr) = 2^{n+1}$，
$s(n) = 2^{n+1} - 2$。 因此 $n \ge 2$ 時
\[
f(n) = s(n-1) + 2 = 2^{n} - 2 + 2 = 2^n, \qquad b(n) = 2^n - 2 .
\]
逐項驗證：$f(2) = 1 + 1 + 2 = 4$，$b(2) = 2$；$f(3) = 4 + 2 + 2 = 8$，$b(3) = 6$；$f(4) = 16$……
foo(10) 回傳 $2^{10} = \mathbf{1024}$。

\textbf{㈢} 每次呼叫除了遞迴，只做常數個運算。 令 $T_f(n)$、$T_b(n)$ 為 foo($n$)、bar($n$) 的時間，$n \ge 2$ 時
$T_f(n) = T_f(n-1) + T_b(n-1) + c$，$T_b(n)$ 也是同一個式子，所以兩者相同，記為 $T(n)$：
\[
T(n) = 2T(n-1) + c, \qquad T(1) = c',
\]
展開（資結書第~2.1~節）：$T(n) = 2^{n-1}c' + (2^{n-1} - 1)c = \Theta(2^n)$。 也可以直接數呼叫次數：foo($n$) 共
呼叫 foo $2^{n-1}$ 次、bar $2^{n-1} - 1$ 次，合計 $2^n - 1$ 次，每次 $O(1)$。 所以 foo($n$) 的時間複雜度是
$O(2^n)$（而且是 $\Theta(2^n)$）。 遞迴深度只有 $n$，堆疊空間 $O(n)$。 若改成記住算過的結果（動態規劃），
每個 $(\text{函式}, n)$ 只算一次，時間降為 $O(n)$——題目特別排除了這種做法。
\end{pastsoln}
